% !TeX root = ../main.tex

\chapter{绪论}

\section{研究背景与意义}

\subsection{电解水制氢与波动性电源适配}

实现碳达峰、碳中和（``双碳''目标）对能源系统的低碳转型提出了迫切需求。
氢能作为零碳排放的二次能源，在化工、冶金、交通等难电气化领域的深度脱碳中具有不可替代的作用。
利用可再生能源电解水制取的``绿氢''，其全生命周期碳排放接近于零，
已成为当前氢能发展的主要技术路线。

目前，电解水制氢技术主要包括碱性电解水（Alkaline Water Electrolysis, AWE）、
质子交换膜电解水（Proton Exchange Membrane, PEM）和固体氧化物电解水（Solid Oxide Electrolysis Cell, SOEC）三种主流技术路线。
其中，AWE技术已有数十年工业化应用历史，单体设备成本低、系统成熟度高，
是目前兆瓦级及以上制氢项目的主流选择~\cite{Hu2022AppliedEnergy}。
PEM电解槽响应速度快、负载范围宽，但设备成本较高且对水质要求严格；
SOEC在高温下运行，理论效率潜力大，但材料热循环稳定性与动态响应能力
仍有待提升~\cite{Zhang2025TEC}。
近年来，ALK与PEM混合制氢系统因兼顾宽功率适应性与高效率而受到关注，
研究表明合理的容量配比可在提升可再生能源利用率的同时降低综合投资成本~\cite{Wang2026APEN}。

\subsection{波动性电源适配挑战}

随着风电、光伏等可再生能源装机规模的快速增长，其间歇性和波动性特征对电网稳定运行构成了持续压力。
电解水制氢系统作为大规模可调节负荷，是消纳波动性可再生能源的重要途径。
但这也要求电解系统具备更强的宽功率运行能力：

\begin{itemize}
  \item \textbf{宽功率范围}：需要在20\%--150\%额定功率范围内稳定运行
  \item \textbf{快速响应}：需跟踪秒级至分钟级的功率波动
  \item \textbf{安全约束}：需确保温度、氢气纯度等关键参数始终处于安全范围
\end{itemize}

上述挑战在风电直接驱动的离网制氢场景中尤为突出。
Zhang等~\cite{Zhang2021IJHE_pilot}在250~kW风电制氢示范项目中验证了
碱性电解槽可在30\%--100\%额定功率范围内动态运行，
但温度上升的分钟级惯性严重限制了可再生能源的利用率；
Haoran等~\cite{Haoran2024IJHE}系统综述了可再生能源驱动碱性电解槽
宽范围运行时的安全与效率问题，指出功率调节灵活性、气体纯度和热一致性
是制约大规模应用的关键瓶颈。

\subsection{多槽协调控制难点}

大型电解制氢项目通常采用多槽并联的模块化架构以提高系统灵活性和可靠性。
四槽并联系统是常见的工业配置，
Qiu等~\cite{Qiu2025arXiv}针对此类N-in-1多槽AWE系统建立了状态空间模型，
揭示了碱液循环、温度与HTO杂质在多槽间的动态耦合机理；
Zeng等~\cite{Zeng2025TSTE}进一步指出，晶闸管整流器供电的多电解槽机组
存在有功-无功非线性耦合，不合理的负荷分配会导致电压越限与网损增加。
综合来看，多槽协调控制面临以下挑战：

\begin{enumerate}
  \item \textbf{热耦合}：多个电解槽共用换热器和分离器，形成复杂的热力学耦合
  \item \textbf{气相耦合}：分离器气相空间共享，氢气传输相互影响
  \item \textbf{功率分配}：总功率约束下各槽的最优负荷分配问题
  \item \textbf{安全约束}：氢气至氧侧传输量（HTO）等关键安全指标的严格约束
  \item \textbf{退化均衡}：各槽运行时间与启停次数差异导致寿命不均衡~\cite{Zhou2025RE}
\end{enumerate}

\section{国内外研究现状}

\subsection{电解槽控制}

在电解槽控制领域，传统方法以比例-积分-微分（PID）反馈控制和基于规则的启发式策略为主。
PID控制器结构简单、易于实现，但其本质为无模型的比例-积分-微分反馈律，
仅依赖当前时刻的观测误差进行校正，既无法利用未来功率计划等预测信息，
也缺乏对多变量耦合约束的最优性保证，因而在处理具有大惯性、强耦合和非线性特征的
电解系统时，动态响应和安全约束满足能力有限~\cite{Li2024IJHE_chemeng}。
为拓宽负荷调节范围，Qi等~\cite{Qi2021IJHE}提出了基于压力调节的控制策略，
通过稳态与动态杂质积累模型设计了最优控制时间（OCT）与模型预测控制（MPC）控制器，
将系统的最小负荷从27.5\%分别降至10.5\%和10\%。

在电-氢耦合系统层面，Xia等~\cite{Xia2025TPEL}研究了离网光伏制氢场景下
电解槽温度变化对功率调节能力的约束，并提出电池预热策略以提升功率消纳灵活性，
使2.5~kW电解槽的最大可消纳功率从额定功率的59.12\%提升至96.56\%；
Zhang等~\cite{Zhang2025EPSR}则基于灰色关联分析量化了影响制氢效率的敏感因素，
提出了改进的电压-功率下垂控制策略，有效提升了风电制氢效率。
此外，Wang等~\cite{Wang2025TSG}面向并网绿氢混合储能系统设计了电压源分层控制框架，
在考虑运行约束和电流变化率限制的同时实现了氢储能与电网间的功率分配，
并通过800~V直流微电网硬件在环实验进行了验证。

随着制氢规模扩大，多槽并联成为主流架构，其协调控制问题近年受到广泛关注。
在系统建模层面，Qiu等~\cite{Qiu2025arXiv}针对N-in-1多槽AWE系统建立了状态空间模型，
描述了碱液循环、温度与HTO杂质在多槽间的动态耦合行为，并采用非线性模型预测控制（NMPC）协调优化各槽电流分配、碱液流量与冷却。
Qiu等~\cite{Qiu2023Renewable}进一步面向数十台电解槽组成的工业级P2H工厂，
建立了多物理场感知的最优生产调度模型（MILP），显式考虑动态热力学与HTO杂质传输过程，
通过分解算法实现了大规模系统的可扩展求解，显著提升了氢气产量与收益。

在协调控制策略方面，研究沿着多个技术路线展开。
基于优化控制的路径中，Yang等~\cite{Yang2025IJHE}分析了波动工况下温度变化对AWE安全边界的影响，
提出了基于可行域分析的参数协调控制方法，有效降低了低负荷工况下的HTO浓度；
Zheng等~\cite{Zheng2022AppliedEnergy}将热-电耦合模型用于日前优化调度，
发现部分负荷运行在经济性上优于满负荷恒定运行。
Li等~\cite{Li2025TPEL}面向大规模并联AWE系统，提出了基于下垂系数自适应调整和模糊PI的鲁棒协调控制策略，
实验验证了其在功率分配准确性和动态响应方面的优势。
Chen等~\cite{Chen2025TSTE}采用均值场控制（MFC）方法处理大规模电解槽阵列的功率分配问题，
通过HJB-FPK方程描述随机功率波动和电解槽故障下的群体行为，实现了快速功率分解。
Wang等~\cite{Wang2025TSTE}进一步考虑多电解槽的寿命退化与利用均衡，
提出了功率-状态滚动优化策略，显著提升了系统效率和氢气平准化成本。

Urs{\'u}a等~\cite{Survey2024Hydrogen}对电解系统控制策略进行了全面综述，
归纳了从PID到先进模型预测控制的多种方法，指出MPC在显式处理约束方面具有理论优势，
但计算开销随系统规模呈指数级增长，难以直接扩展至大规模多槽系统。

综合上述研究可见，现有电解槽控制方法主要依赖传统优化或反馈控制理论：
NMPC和MILP需在每个控制周期内在线求解复杂优化问题，单次求解耗时数秒至数分钟，
在分钟级控制周期下难以满足实时性需求；
而基于规则或线性化的控制策略（如下垂控制、模糊PI）虽计算高效，
但难以充分挖掘多槽系统在多物理场耦合条件下的最优运行潜力。
这促使研究者从两个方向寻求突破：一是提升控制策略的时间分辨率与集群可扩展性；二是在优化目标中显式纳入设备寿命与可再生能源波动特性。

在提升时间分辨率与集群可扩展性方面，
Guan等~\cite{Guan2025TSTE}面向风电-氢能耦合场景，建立了涵盖
电化学、传热与杂质传输的精细化动态模型，并据此设计了分钟级优化运行策略。
通过碱液流量、冷却流量和压力的协同调节，该策略实现了1分钟间隔的功率跟踪，
将电解槽负荷范围拓宽13.8\%，弃风率降低15.06\%。
Wen等~\cite{Wen2026ECM}提出了一种通用的电解槽集群实时控制框架，
通过规则-求解器混合优化策略显式考虑电流效率、温度、性能退化和设备异质性，
在案例研究中实现了氢产量最大提升9.27\%、系统效率提升5.89\%。
在设备退化均衡方面，Zhou等~\cite{Zhou2025RE}提出了考虑电解槽退化的
多电解槽集群分配策略，将温度、HTO和OTH杂质变化纳入功率分配框架，
有效降低了各槽退化差异；
Lei等~\cite{Lei2026APEN}则进一步设计了多目标排序轮换控制策略，
通过Pareto前沿均衡多槽磨损，日常运行中退化标准差较传统启停策略降低88\%。

在可再生能源耦合与电网互动方面，
Zhang等~\cite{Zhang2025IJHE_solar}在光伏制氢领域构建了多碱性电解槽
协同优化框架，融合长短期记忆网络-卷积神经网络（LSTM-CNN）功率预测与粒子群动态调度算法，
将光伏利用率提升至98.6\%以上。
在电网互动与能量管理方面，
Zeng等~\cite{Zeng2025TSTE}注意到，晶闸管整流器供电的多电解槽机组
存在有功-无功非线性耦合，进而建立了有功-无功协调管理方法，
通过混合整数二阶锥规划优化负荷分配与热特性，
在提升氢气产量的同时降低网损。
Li等~\cite{Li2024APEN_two_layer}则面向并网多槽制氢站提出了双层能量管理策略，
上层基于模型的鲁棒调度优化功率因数与生产目标，
下层采用规则算法进行实时增量修正，
在保障控制精度的同时兼顾电网侧功率因数约束。
在多能源系统协同优化方面，Huang等~\cite{Huang2022ECM}将兆瓦级
碱性电解槽的氢-热-电动态与余热回收纳入经济模型预测控制框架，
相较传统规则策略实现运行成本降低59\%。
上述研究从调度、控制与能量管理多角度推进了多槽AWE系统的协调优化，
但仍主要依赖在线求解复杂优化问题，实时性瓶颈尚未根本解决。
相较而言，本文同样以功率分配与温度调节为核心控制目标，
但放弃在线优化路径，转而采用基于扩散模型的生成式策略逼近NMPC专家行为，
在保持较高控制精度的同时将单次推理时间从数秒降至约50毫秒，
从根本上突破了NMPC类方法的实时性瓶颈。

\subsection{电解槽建模}

AWE系统的建模研究涵盖电化学、热力学和气体纯度三个子模型，其建模精度直接决定了后续控制策略的有效性。
Hu等~\cite{Hu2022AppliedEnergy}对现有AWE数学建模工作进行了系统综述，
指出电化学建模研究相对成熟，但热力学建模与气体纯度建模的研究仍较少，
模型预测精度和基于模型的控制策略有待进一步发展。
Huang等~\cite{Huang2025RSER}进一步从系统级决策视角对碱性电解槽建模进行了分类梳理，
将现有模型分为线性电-氢模型（LEHM）、非线性电-氢模型（NEHM）和综合电热氢模型（IEHHM）三类，
指出IEHHM因显式集成了热力学动态而更适合波动工况下的优化决策，但计算复杂度也相应增加。

在动态工况建模与实验验证方面，Gu等~\cite{Gu2024JPS}基于250 kW工业级AWE系统开展了启停和变负荷多时间尺度动态实验，
揭示了HTO响应对参数变化存在分钟级时间延迟和惯性环节，温度是启停过程的主要限制因素。
Zheng等~\cite{Zheng2022AppliedEnergy}建立了涵盖氢气产率、温度变化与状态切换（运行/待机/停机）的综合电解槽模型，
用于日前优化调度，发现该模型能有效捕捉热-电动态耦合，但计算开销显著高于传统简化模型。
Qi等~\cite{Qi2022AppliedEnergy}建立了碱性电解系统的热力学模型并设计了温度控制器；
Zhong等~\cite{Zhong2025AppliedEnergy}进一步针对波动可再生能源输入下的非稳态热力学特性进行了分析与优化。
然而，上述工作主要针对单槽系统，对于多槽共用辅助设备（如气液分离器、碱液循环回路）带来的耦合动力学问题，建模研究尚不充分。

在实验建模与多物理场仿真方面，Ren等~\cite{Ren2022JPS}对250~kW工业级
碱性电解槽开展了系统实验与半经验建模，揭示了冷启动至额定工况需6小时以上，
且运行条件对效率与安全具有显著影响；该模型在风电制氢示范项目中得到了验证。
Yin等~\cite{Yin2025RE}基于Modelica语言建立了100~kW碱性电解系统的多物理场模型，
涵盖电解槽、换热器、气液分离器等关键设备，
通过稳态与动态仿真分析了热惯性和容积惯性对温度和产气响应时间的主导作用。
Shi等~\cite{Shi2023IECR}则从全流程视角出发，建立了工业碱性电解水厂的
plant-wide模型，并设计了考虑关键变量积分特性的分层实用控制策略。
在氢气纯度建模方面，Haug等~\cite{Haug2017IJHE}通过在线气相色谱实验，
系统考察了温度、电解液浓度和流量对零间隙碱性电解槽产品气纯度的影响，
发现升高温度和降低流量可减少气体杂质，
为动态工况下的纯度控制提供了实验依据。

综合来看，现有AWE系统建模研究多针对单槽系统或简化动态模型，
对多槽并联架构下共用辅助设备（如气液分离器、碱液循环回路）带来的耦合动力学刻画尚不充分。
本文则面向四槽并联系统，建立了涵盖电化学、热力学与氢气传输的高保真耦合动力学模型，
显式刻画了碱液循环、热耦合、气相耦合与功率分配机理，
为后续控制策略设计提供了更为精细的物理基础。

\subsection{基于学习的控制方法}

人工智能技术为工业控制提供了新的解决思路。
强化学习（RL）通过与环境交互学习最优策略，已在游戏控制和机器人操作中取得成功~\cite{Mnih2015Nature}。
然而，RL需要大量探索性试验以收集经验数据，电解系统作为高安全要求的化工过程，
不允许在真实设备上进行危险的随机探索，这限制了RL的直接应用。

模仿学习（Imitation Learning）通过模仿专家行为学习策略，避免了危险的探索过程~\cite{Ross2011AISTATS}。
其中，行为克隆（Behavior Cloning）通过监督学习建立从状态到动作的确定性映射，
实现最为简单，但存在两个固有缺陷：
一是分布漂移（Distribution Shift）——训练时从专家状态分布采样，而推理时从策略自身产生的状态分布采样，两者不一致导致误差累积；
二是单峰输出假设——高斯分布只能输出单一均值，难以刻画多槽功率分配问题中存在的多种等价最优负荷分配方案。

在固体氧化物电解槽（SOEC）控制方面，
Zhang等~\cite{Zhang2025TEC}采用双延迟深度确定性策略梯度（TD3）算法
实现了SOEC制氢系统的需求响应最优控制，
在跟踪电网调度信号的同时保持高能量效率。
扩散模型（Diffusion Models）通过学习逆转一个逐步加噪的扩散过程来生成高质量数据，
在图像生成、语音合成等领域表现突出，近年来逐渐应用于控制任务。
在机器人控制领域，Chi等~\cite{Chi2023DiffusionPolicy}提出的扩散策略（Diffusion Policy）
通过建模动作的条件分布而非确定性映射，有效解决了行为克隆的分布漂移和单模态输出问题，
在复杂操作技能学习中展现出丰富的动作分布建模能力。

在控制与规划任务中，扩散模型与模型预测控制（MPC）的结合逐渐受到关注。
Julbe等~\cite{Julbe2025AMPC}提出基于扩散模型的近似MPC方法，
利用扩散模型完整表示MPC解的分布，在7自由度机械臂控制任务中实现了超过70倍的在线求解加速，
同时通过代价引导机制在闭环中保持模态一致性。
R{\"o}mer等~\cite{Roemer2025DPCC}提出带约束的扩散预测控制（DPCC），
在去噪过程中嵌入基于模型的投影和约束收紧机制，使扩散策略能够满足训练数据中未见过的新型约束，
在机械臂规划任务中验证了其对未见约束的泛化能力。

然而，上述数据驱动与生成模型方法主要针对单槽系统或简化的动态模型，
对于多槽并联系统中热-气-电多物理场耦合带来的高维状态空间与复杂安全约束，
其可扩展性和安全保证仍面临挑战。
相较而言，本文放弃高斯单峰假设与确定性映射，
采用扩散策略（Diffusion Policy）建模动作的条件分布，
有效解决了行为克隆中的分布漂移与单模态输出问题；
同时通过代价加权动作选择机制在候选序列中优选低代价动作，
在保持较高控制精度的同时将单次推理时间降至约50毫秒，
为电解槽等高维耦合系统的实时控制提供了新的技术路径。

为明确各类方法在本研究中的定位与适用边界，从核心原理、优势与局限性三个维度对PID、NMPC与扩散策略进行对比。对比结果如表~\ref{tab:method_comparison} 所示。

\begin{table}[!htbp]
  \centering
  \caption{不同控制方法对比}
  \label{tab:method_comparison}
  \begin{tabular}{@{}l p{2.8cm} p{3.2cm} p{3.8cm}@{}}
    \toprule
    \textbf{方法} & \textbf{核心原理} & \textbf{主要优势} & \textbf{局限性} \\
    \midrule
    PID & 比例-积分-微分反馈 & 结构简单，易于实现，鲁棒性好 & 无最优性保证，仅利用当前观测信息，难以处理强耦合、大惯性与严格约束 \\
    \addlinespace
    NMPC & 在线非线性优化 & 理论完备，显式处理约束 & 计算开销大，实时性差 \\
    \addlinespace
    Diffusion Policy & 去噪扩散生成 & 高质量采样 & 推理步数多，耗时较长 \\
    \bottomrule
  \end{tabular}
\end{table}

\subsection{安全控制的方法}

在电解槽等高安全要求的化工过程中，控制策略必须严格保证温度、气体纯度等关键指标不越界。
Su等~\cite{Su2026IJHE}针对碱性电解槽的宽功率波动工况，提出了基于安全强化学习的自抗扰控制策略，
通过设计安全层约束探索过程，避免了危险状态下的随机试错，
在仿真中验证了该方法对温度越限和气体纯度异常的抑制效果。
然而，安全强化学习本质上仍依赖值函数估计与策略梯度更新，
对温度、HTO等时变非线性约束缺乏严格的数学保证。

控制障碍函数（Control Barrier Functions, CBF）为安全约束提供了严格的前向不变性理论框架。
Ames等~\cite{Ames2019CBF}系统阐述了CBF在仿射控制系统中的设计方法，
通过构造满足李导数条件的障碍函数，确保系统状态始终停留在安全集合内部。
近年来，CBF在机器人避障与无人机安全控制中得到广泛应用，
但在具有非仿射非线性特性的化工过程控制中，CBF的显式构造与在线求解仍面临挑战。

安全投影算子与安全滤波是另一类重要的安全增强手段，
通过在策略输出端引入修正层，将潜在不安全的控制量映射到可行控制集。
R{\"o}mer等~\cite{Roemer2025PACS}提出了路径一致的安全滤波（Path-Consistent Safety Filtering）框架，
通过对扩散策略生成的轨迹进行路径一致的制动修正，在保持任务成功率的同时提供形式化安全保证；
同一团队~\cite{Roemer2025DPCC}进一步将基于模型的投影操作嵌入扩散去噪迭代，
实现了对训练和测试时约束的显式满足。
上述方法在机器人路径规划中的凸约束集合上表现良好，
但对于电解系统中多物理场耦合、时变非线性的复杂安全约束，其可扩展性与严格保证能力仍显不足。

在扩散模型与安全约束的结合方面，
Xiao等~\cite{Xiao2023SafeDiffuser}提出了SafeDiffuser框架，
通过在扩散去噪过程中引入安全值函数引导，使生成轨迹满足碰撞避免等简单几何约束。
然而，现有扩散模型安全增强研究主要面向机器人规划中的已知凸约束集合，
对于工业过程中多物理场耦合、时变非线性的复杂安全约束，其可扩展性与严格保证能力仍显不足。

相较而言，本文针对生成模型输出存在尾部风险的问题，
结合投影算子单步预测修正与CBF李导数前向不变性理论，
构建了适用于多槽AWE系统的安全控制框架。
投影算子通过非线性规划将无约束策略输出映射到可行控制集；
CBF推导则显式处理了电流与法拉第效率之间的非仿射耦合项，并给出了严格的理论保证。

\subsection{研究空白与本文切入点}

综合上述分析，现有研究在AWE系统协调控制中存在以下空白：
\begin{enumerate}
  \item \textbf{多槽耦合动力学刻画不足}：现有AWE系统建模研究多针对单槽系统或简化动态模型，
  对多槽并联架构下共用辅助设备带来的热-气-电耦合动力学刻画尚不充分，
  难以为后续协调控制策略提供精细的物理基础。
  \item \textbf{计算效率瓶颈}：NMPC和MILP等基于在线优化的方法单次求解耗时数秒至数分钟，
  在分钟级控制周期下难以满足实时性需求~\cite{Survey2024Hydrogen, Guan2025TSTE}；
  而基于规则或线性化的策略~\cite{Li2024IJHE_chemeng, Zhang2025EPSR}虽响应迅速，
  却难以充分挖掘多槽耦合条件下的最优运行潜力。
  \item \textbf{动作分布建模不足}：传统模仿学习与行为克隆方法采用高斯单峰假设，
  难以刻画多槽系统功率分配中存在的多种等价最优解~\cite{Ross2011AISTATS}；
  强化学习虽具探索能力，但在高安全要求的化工过程中难以开展危险随机试验~\cite{Mnih2015Nature, Su2026IJHE}。
  \item \textbf{安全约束保证缺失}：数据驱动方法（包括生成模型）在学习与推理过程中
  缺乏对温度、HTO等关键安全约束的严格保证~\cite{Xiao2023SafeDiffuser}。
  现有安全滤波研究主要针对机器人规划中的已知约束集合，
  对于工业过程中多物理场耦合、时变非线性的复杂安全约束，
  如何设计显式的安全修正机制仍是开放性问题。
\end{enumerate}

为填补上述研究空白，本文面向四槽并联AWE系统的功率分配与温度调节问题，
建立基于扩散模型的安全预测控制框架，
在保持NMPC级控制精度的同时将推理时间降至毫秒级，
并严格保证温度、氢气纯度等安全约束的前向不变性。

\section{研究内容与主要贡献}

\subsection{研究内容}

围绕四槽并联AWE系统的协调控制问题，本文从物理建模、策略学习、安全保证和实验验证四个层面展开研究：

\begin{enumerate}
  \item \textbf{高保真物理建模}：建立涵盖电化学、热力学和氢气传输的多槽系统耦合动力学模型，分析各槽间的热耦合、气相耦合与功率耦合机理；

  \item \textbf{专家数据生成}：设计NMPC基线控制器，在典型风电功率曲线下生成覆盖多工况的专家状态-动作数据集；

  \item \textbf{生成式控制策略}：提出基于Diffusion TCN的预测控制策略，利用扩散模型的生成能力生成候选动作序列，并通过代价加权选择机制优化控制品质；

  \item \textbf{安全约束保证}：引入投影算子与CBF安全滤波两类约束保证机制，将生成模型的无约束输出修正为严格满足安全边界的可行动作；

  \item \textbf{实验验证}：构建完整的闭环仿真平台，在功率跟踪、温度调节、安全约束满足率和计算效率四个维度上与NMPC基线进行对比验证。
\end{enumerate}

\subsection{主要贡献}

本研究主要贡献包括：

\begin{itemize}
  \item \textbf{将扩散模型引入AWE系统协调控制}：
  针对多槽电解系统的功率分配与温度调节问题，提出基于Diffusion TCN的预测控制策略，
  利用生成模型的采样能力捕捉多种等价最优动作模式，并通过代价加权选择机制实现动作优选，
  将单次推理时间从NMPC的2--5秒降低至约50毫秒。

  \item \textbf{设计了面向多槽系统的安全约束保证机制}：
  针对生成模型输出存在尾部风险的固有缺陷，结合投影算子单步预测修正与CBF李导数前向不变性理论，
  构建了适用于多槽AWE系统的安全控制框架。其中CBF推导显式处理了电流-法拉第效率的非仿射耦合项，
  并给出了严格的理论保证。

  \item \textbf{建立了完整的仿真-训练-验证实验平台}：
  开发了四槽AWE系统的高保真物理仿真器，涵盖电化学、传热与氢气传输的多时间尺度耦合动力学；
  实现了从NMPC专家数据生成、扩散模型训练到安全闭环控制的完整实验流程，
  在正常工况与高温、高功率、低功率HTO累积等极端工况下验证了所提方法的有效性与实时性。
\end{itemize}

\section{论文组织结构}

全文共分为四章，各章之间的逻辑关系如下：

\textbf{第一章 绪论}阐述AWE系统协调控制的研究背景，梳理传统控制、AI驱动控制和生成模型方法的研究现状，并指出当前研究在实时性与安全性方面的不足，明确研究内容与贡献。

\textbf{第二章 AWE系统建模与仿真}首先建立单槽系统的电化学、热力学与氢气传输模型，在此基础上引入多槽并联架构与共用辅助设备带来的耦合机制，逐步扩展得到多槽耦合系统的统一状态空间模型，并明确控制问题中的状态向量、控制输入与控制目标，最终搭建完整的仿真平台。

\textbf{第三章 基于扩散模型的安全学习控制框架}是核心章节。该章对比不同控制方法的技术路线，阐述所提``专家数据生成--扩散模型训练--安全约束闭环''整体框架，提出基于Diffusion TCN的生成式控制策略，给出投影算子和控制障碍函数（CBF）的安全约束保证机制及完整数学推导，并通过闭环仿真验证所提方法的有效性。

\textbf{第四章 结论与展望}总结主要结论与创新点，分析研究的局限性，并展望未来可能的改进方向。
